\section{Related Work}
\label{sec:related}

\paragraph{Trained interpretability probes.} Sparse Autoencoders \citep{bricken2023towards,cunningham2023sparse} decompose activations into learned dictionary atoms with L1 or TopK sparsity. \vpd\ \citep{bushnaq2026vpd} extends this to weight decomposition with stochastic masking plus a learned importance gate; \bsf\ \citep{goodfire2026bsf} generalizes to block-level sparsity via Grassmannian, block, and group-lasso variants, and reports 2--4-dimensional concept structure on vision. Concurrent Anthropic work on Jacobian lenses \citep{anthropic2026jlens} computes a linearized projection of activations onto vocabulary directions, sharing our observation that analytic linearizations often approximate trained probes' outputs. All of these methods pay per-layer training cost; \svdomp\ removes that cost while matching or exceeding their weight-decomposition quality.

\paragraph{SVD and compressed sensing.} The Eckart-Young theorem \citep{eckart1936approximation} establishes Frobenius optimality of top-$k$ SVD; Weyl and Davis-Kahan perturbation bounds \citep{davis1970rotation} give quantitative stability. Orthogonal matching pursuit \citep{pati1993omp} in the compressed-sensing tradition gives approximate sparse recovery guarantees for general dictionaries; when the dictionary is orthogonal, OMP collapses to top-$k$ closed form. Our contribution is not a new sparse-coding method but the observation that these classical results already dominate trained-probe methods on the specific task of weight decomposition for interpretability.

\paragraph{Efficient interpretability tooling.} Recent work on distillation of interpretability methods \citep{efficient2026distill} and low-rank approximations to SAE dictionaries \citep{low2026sae} aims to reduce inference cost of trained probes. \svdomp\ is orthogonal: it eliminates the underlying training cost rather than distilling a trained model.
